% ৪.৫ ওয়েবসাইটের প্রকারভেদ: Static, Dynamic, Local, Remote
\section{ওয়েবসাইটের প্রকারভেদ}

\begin{learningbox}
এই টপিক শেষে শিক্ষার্থী Static ও Dynamic ওয়েবসাইটের মধ্যে পার্থক্য, এবং Local ও Remote ওয়েবসাইট কী তা ব্যাখ্যা করতে পারবে।
\end{learningbox}

\subsection{ভূমিকা}
ওয়েবসাইট সবসময় একই ধরনের হয় না। কোনো ওয়েবসাইটে তথ্য স্থির থাকে; আবার কোনো ওয়েবসাইটে ব্যবহারকারীর input বা database-এর ওপর ভিত্তি করে তথ্য পরিবর্তিত হয়। অবস্থানের ভিত্তিতেও website local বা remote হতে পারে। এই টপিকটি HSC ICT-এ প্রাসঙ্গিক কনসেপ্টগুলো সরলভাবে তুলে ধরে।

\subsection{Static Website}
\begin{definitionbox}{Static Website}
Static Website হলো এমন ওয়েবসাইট যার content সাধারণত fixed থাকে এবং ব্যবহারকারীভেদে স্বয়ংক্রিয়ভাবে পরিবর্তিত হয় না।
\end{definitionbox}

\begin{keypointbox}{Static ওয়েবসাইটের বৈশিষ্ট্য}
\begin{itemize}
\item Content fixed থাকে; সাধারণত HTML/CSS দিয়ে তৈরি।
\item Database connection থাকে না।
\item দ্রুত load হয়; তৈরি ও maintain করা সহজ।
\item উপযোগী: portfolio, brochure, notice page ইত্যাদি।
\end{itemize}
\end{keypointbox}

\subsection{Dynamic Website}
\begin{definitionbox}{Dynamic Website}
Dynamic Website হলো এমন ওয়েবসাইট যার content user input, database বা অন্যান্য condition অনুযায়ী পরিবর্তিত হতে পারে।
\end{definitionbox}

\begin{keypointbox}{Dynamic ওয়েবসাইটের বৈশিষ্ট্য}
\begin{itemize}
\item Content পরিবর্তনশীল; database ও server-side scripting ব্যবহার করে।
\item Login, registration, comment, search, online order ইত্যাদি সম্ভব।
\item উপযোগী: e-commerce, social media, online result system ইত্যাদি।
\end{itemize}
\end{keypointbox}

\subsection{Static vs Dynamic (তালিকা)}
\begin{center}
\begin{tabular}{|p{0.22\textwidth}|p{0.29\textwidth}|p{0.29\textwidth}|}
\hline
\textbf{বিষয়} & \textbf{Static} & \textbf{Dynamic} \\
\hline
Content & Fixed & পরিবর্তনশীল \\
\hline
Database & না & আছে \\
\hline
User interaction & কম & বেশি \\
\hline
Load speed & দ্রুত & তুলনামূলক ধীর \\
\hline
Use cases & Portfolio, brochure & E-commerce, portals \\
\hline
\end{tabular}
\end{center}

\subsection{Local Website}
\begin{definitionbox}{Local Website}
Local Website হলো এমন website বা web page যা নিজের computer বা local server-এ থাকে এবং সাধারণত development/testing-এর জন্য ব্যবহৃত হয়।
\end{definitionbox}

\begin{keypointbox}{Local ওয়েবসাইটের বৈশিষ্ট্য}
\begin{itemize}
\item Offline বা local environment-এ test করা যায় (XAMPP, WAMP, Live Server ইত্যাদি)।
\item Hosting cost লাগে না; ত্রুটি পরীক্ষা ও rapid development-এর জন্য উপযোগী।
\end{itemize}
\end{keypointbox}

\subsection{Remote Website}
\begin{definitionbox}{Remote Website}
Remote Website হলো এমন website যা দূরবর্তী web server/hosting server-এ রাখা থাকে এবং internet-এর মাধ্যমে public access করা যায়।
\end{definitionbox}

\begin{keypointbox}{Remote ওয়েবসাইটের বৈশিষ্ট্য}
\begin{itemize}
\item অনলাইনে public user access করতে পারে; domain ও hosting প্রয়োজন।
\item Maintenance, backup ও security জরুরি।
\end{itemize}
\end{keypointbox}

\subsection{Local vs Remote (তালিকা)}
\begin{center}
\begin{tabular}{|p{0.20\textwidth}|p{0.30\textwidth}|p{0.30\textwidth}|}
\hline
\textbf{বিষয়} & \textbf{Local} & \textbf{Remote} \\
\hline
অবস্থান & নিজের computer/local server & hosting/remote server \\
\hline
Access & নিজের device থেকে & world-wide via internet \\
\hline
Use case & development/testing & production/public use \\
\hline
Cost & কম/নেই & hosting/domain cost লাগে \\
\hline
\end{tabular}
\end{center}

\begin{examplebox}{Website type diagram}
Static: User → Browser → Web Server → Fixed HTML Page \\
Dynamic: User → Browser → Web Server → Server Script → Database → Generated Page \\
Local vs Remote: Computer → Browser → (Local file) / (Internet → Web Server)
\end{examplebox}

\subsection{Static Website কোথায় ব্যবহার করা ভালো?}

যেসব website-এ তথ্য ঘন ঘন পরিবর্তন হয় না এবং user login/input দরকার হয় না, সেখানে static website যথেষ্ট। HSC practical-এ simple HTML page, college notice page, personal profile page, portfolio page ইত্যাদি static website হিসেবে দেখানো যায়।

\begin{examplebox}{Static website example}
\begin{itemize}
\item Personal portfolio
\item College notice page
\item Product information page
\item Event announcement page
\item Basic company profile
\end{itemize}
\end{examplebox}

\subsection{Dynamic Website কোথায় ব্যবহার করা ভালো?}

যেসব website user input, login, search, database, online payment বা real-time content update ব্যবহার করে, সেগুলো dynamic website। এই ধরনের website সাধারণত HTML/CSS-এর পাশাপাশি server-side language, database ও security system ব্যবহার করে।

\begin{examplebox}{Dynamic website example}
\begin{itemize}
\item Online admission system
\item Online result system
\item E-commerce website
\item Social media platform
\item News portal
\item Online banking বা ticket booking system
\end{itemize}
\end{examplebox}

\subsection{Static থেকে Dynamic হওয়ার উদাহরণ}

একটি college website প্রথমে শুধু notice, image ও contact information দেখালে সেটি static হতে পারে। কিন্তু একই website-এ যদি student login, online result search, admission form submission বা database থেকে notice update করার system থাকে, তাহলে সেটি dynamic website হবে।

\begin{keypointbox}{সহজ সিদ্ধান্ত নেওয়ার নিয়ম}
\begin{itemize}
\item শুধু fixed information দেখালে: Static website।
\item user input নিয়ে result/search/login দেখালে: Dynamic website।
\item নিজের computer-এ test করলে: Local website।
\item internet server-এ upload করে public করলে: Remote website।
\end{itemize}
\end{keypointbox}

\subsection{Website Type Selection Tree}

\begin{center}
\begin{tikzpicture}[
  level distance=1.15cm,
  level 1/.style={sibling distance=4.2cm},
  level 2/.style={sibling distance=2.2cm},
  every node/.style={draw, rounded corners, align=center, fill=yellow!8, inner sep=4pt, font=\small}
]
\node {Website কোন ধরনের?}
  child { node {Content fixed?}
    child { node {হ্যাঁ\\Static} }
    child { node {না\\Dynamic} }
  }
  child { node {কোথায় রাখা?}
    child { node {নিজ computer\\Local} }
    child { node {Internet server\\Remote} }
  };
\end{tikzpicture}
\end{center}

\subsection{Problem Solution}

\begin{practicebox}
\textbf{Problem:} একটি college website-এ শুধু college পরিচিতি, image, subject list এবং contact address আছে। কিন্তু আরেকটি website-এ student roll দিলে database থেকে result দেখা যায়। কোনটি static এবং কোনটি dynamic?

\textbf{Solution:}
\begin{itemize}
\item প্রথম website-এ শুধু fixed information আছে, তাই এটি static website।
\item দ্বিতীয় website user input নেয় এবং database থেকে result দেখায়, তাই এটি dynamic website।
\item যদি প্রথম website নিজের computer-এ test করা হয়, তাহলে সেটি local website।
\item যদি server-এ upload করে internet-এ দেখা যায়, তাহলে সেটি remote website।
\end{itemize}
\end{practicebox}

\begin{boardbox}{Board focus}
Static vs Dynamic এবং Local vs Remote website-এর পার্থক্য HSC exam-এ short question, MCQ ও CQ reasoning হিসেবে আসতে পারে। উদাহরণসহ উত্তর লিখলে নম্বর ভালো পাওয়া যায়।
\end{boardbox}

\begin{mistakebox}{সাধারণ ভুল}
\begin{itemize}
\item শুধু online আছে বলেই website-কে dynamic বলা।
\item database না থাকলেও সব website-কে dynamic ভাবা।
\item local website মানে ছোট website মনে করা।
\item remote website মানে শুধু বিদেশি website মনে করা।
\item static website-এ কখনো image/link/table থাকতে পারে না--এমন ভুল ধারণা রাখা।
\end{itemize}
\end{mistakebox}


\begin{practicebox}
\begin{enumerate}
\item একটি static ও একটি dynamic website-এর উদাহরণ লেখ।
\item Local ও Remote website-এর মধ্যে ৩টি পার্থক্য লেখ।
\item কেন dynamic website-এ security বেশি গুরুত্বপূর্ণ—সংক্ষিপ্ত ব্যাখ্যা কর।
\item Online result system static না dynamic? কারণসহ লেখ।
\item নিজের computer-এ \texttt{index.html} open করলে সেটি local না remote? ব্যাখ্যা কর।
\end{enumerate}
\end{practicebox}
